\usetheme{Madrid}
\usecolortheme{seahorse}
\setbeamertemplate{navigation symbols}{}
\setbeamertemplate{itemize items}[circle]
\setbeamertemplate{enumerate items}[default]

\usepackage{amsmath,amssymb,mathtools}

\newcommand{\R}{\mathbb{R}}
\newcommand{\N}{\mathbb{N}}
\newcommand{\Prob}{\mathbb{P}}
\newcommand{\E}{\mathbb{E}}
\newcommand{\A}{\mathcal{A}}
\newcommand{\B}{\mathcal{B}}
\newcommand{\F}{\mathcal{F}}
\newcommand{\one}{\mathbf{1}}
\newcommand{\Leb}{\lambda}
\newcommand{\vikiname}[1]{\par\smallskip{\scriptsize\color{blue!60!black}\textbf{LLMviki:} \texttt{\detokenize{#1}}}}
\newcommand{\blankproof}{%
  \vspace{2mm}
  {\color{gray!55}\rule{\linewidth}{0.4pt}}\par\vspace{5mm}
  {\color{gray!55}\rule{\linewidth}{0.4pt}}\par\vspace{5mm}
  {\color{gray!55}\rule{\linewidth}{0.4pt}}\par\vspace{5mm}
  {\color{gray!55}\rule{\linewidth}{0.4pt}}\par
}
\newcommand{\longblankproof}{%
  \vspace{1mm}
  {\color{gray!55}\rule{\linewidth}{0.4pt}}\par\vspace{4mm}
  {\color{gray!55}\rule{\linewidth}{0.4pt}}\par\vspace{4mm}
  {\color{gray!55}\rule{\linewidth}{0.4pt}}\par\vspace{4mm}
  {\color{gray!55}\rule{\linewidth}{0.4pt}}\par\vspace{4mm}
  {\color{gray!55}\rule{\linewidth}{0.4pt}}\par\vspace{4mm}
  {\color{gray!55}\rule{\linewidth}{0.4pt}}\par
}
\newcommand{\proofblock}[1]{\ifteacher #1 \else \blankproof \fi}
\newcommand{\longproofblock}[1]{\ifteacher #1 \else \longblankproof \fi}

\title[Chapter 1.7]{Chapter 1.7 Product Measures, Fubini's Theorem}
\subtitle{Rick Durrett Probability: Theory and Examples}
\author{Hu Jingwu}
\date{}

\begin{document}

\begin{frame}
  \titlepage
  \vspace{-5mm}
  \begin{center}
    \small Built from the local Chapter 1 LLMviki flash cards.
  \end{center}
\end{frame}

\begin{frame}{Roadmap for 1.7}
  \begin{block}{Learning goal}
    Build product measures and use Tonelli--Fubini as the main engine for switching order of integration.
  \end{block}
  \begin{enumerate}
    \item Product measures on product sigma-fields.
    \item Measurable sections and section functions.
    \item Tonelli for nonnegative functions; Fubini for integrable functions.
    \item Product measures as the measure-theoretic preview of independence.
  \end{enumerate}
\end{frame}

\begin{frame}{Product Measures}
  \begin{block}{Theorem}
    Let $(X,\A,\mu)$ and $(Y,\B,\nu)$ be sigma-finite measure spaces. There is a unique measure
    \[
      \mu\times\nu \quad \text{on } \A\otimes\B
    \]
    such that
    \[
      (\mu\times\nu)(A\times B)=\mu(A)\nu(B),
      \qquad A\in\A,\ B\in\B.
    \]
  \end{block}
  \begin{alertblock}{Pitfall}
    Sets in $\A\otimes\B$ are usually much more complicated than finite unions of rectangles.
  \end{alertblock}
  \vikiname{Product measures; ID: durrett_ch1_product_measure}
\end{frame}

\begin{frame}{Product Measures: Proof Idea}
  \proofblock{
  Start with measurable rectangles and define
  \[
    m(A\times B)=\mu(A)\nu(B).
  \]
  Extend this to the algebra generated by rectangles and verify countable additivity there, using sigma-finiteness to localize to finite measure pieces. The extension theorem then produces a measure on $\A\otimes\B$. Uniqueness follows because two such measures agree on a pi-system of rectangles that generates $\A\otimes\B$.
  }
  \vikiname{Product measures; ID: durrett_ch1_product_measure}
\end{frame}

\begin{frame}{Measurable Sections}
  \begin{block}{Theorem}
    If $C\in\A\otimes\B$, define sections
    \[
      C_x=\{y:(x,y)\in C\},\qquad C^y=\{x:(x,y)\in C\}.
    \]
    Then $C_x\in\B$ and $C^y\in\A$. Moreover, if $\nu$ is sigma-finite, then
    \[
      x\mapsto \nu(C_x)
    \]
    is $\A$-measurable.
  \end{block}
  \vikiname{Measurable sections; ID: durrett_ch1_sections_measurable}
\end{frame}

\begin{frame}{Measurable Sections: How to Prove}
  \proofblock{
  Verify the statement first for rectangles $C=A\times B$, where $C_x=B$ if $x\in A$ and $C_x=\varnothing$ otherwise. Then consider the class of sets for which the conclusion holds. This class contains rectangles and is closed under complements and countable disjoint unions. For section measures, prove first for finite-measure rectangles and then use monotone class arguments plus sigma-finite approximation.
  }
  \vikiname{Measurable sections; ID: durrett_ch1_sections_measurable}
\end{frame}

\begin{frame}{Tonelli and Fubini}
  \begin{block}{Tonelli}
    If $f\ge0$ is measurable on $X\times Y$, then
    \[
      \int f\,d(\mu\times\nu)
      =
      \int_X\int_Y f(x,y)\,\nu(dy)\,\mu(dx)
      =
      \int_Y\int_X f(x,y)\,\mu(dx)\,\nu(dy),
    \]
    allowing the value $+\infty$.
  \end{block}
  \begin{block}{Fubini}
    If $\int |f|\,d(\mu\times\nu)<\infty$, the same identities hold with finite signed integrals.
  \end{block}
  \vikiname{Fubini and Tonelli theorem; ID: durrett_ch1_fubini_tonelli}
\end{frame}

\begin{frame}{Tonelli and Fubini: How to Prove}
  \proofblock{
  For indicators $f=\one_C$, the result is exactly the measurable-section theorem:
  \[
    \int \one_C\,d(\mu\times\nu)=\int_X \nu(C_x)\,\mu(dx).
  \]
  Extend by linearity to nonnegative simple functions. For general $f\ge0$, choose simple functions $f_n\uparrow f$ and apply the monotone convergence theorem. For signed $f$, apply Tonelli to $f^+$ and $f^-$ when $\int |f|<\infty$.
  }
  \vikiname{Fubini and Tonelli theorem; ID: durrett_ch1_fubini_tonelli}
\end{frame}

\begin{frame}{Choosing Tonelli or Fubini}
  \begin{block}{Checklist}
    \begin{itemize}
      \item If $f\ge0$, use Tonelli first; infinite answers are allowed.
      \item If $f$ has signs, prove $\int |f|<\infty$ before switching orders.
      \item To prove integrability, Tonelli is often applied to $|f|$.
      \item To compute measures of regions, write the region by sections and integrate their lengths or masses.
    \end{itemize}
  \end{block}
  \vikiname{Fubini and Tonelli theorem; ID: durrett_ch1_fubini_tonelli}
\end{frame}

\begin{frame}{Preview: Independence}
  \begin{block}{Idea}
    If a joint law factors as a product measure,
    \[
      \mathcal L(X,Y)=\mathcal L(X)\times\mathcal L(Y),
    \]
    then the two coordinates behave independently.
  \end{block}
  \begin{alertblock}{Pitfall}
    Independence is a statement about the joint law, not merely about knowing the two marginal laws.
  \end{alertblock}
  \vikiname{Product measures as preview of independence; ID: durrett_ch1_independence_preview}
\end{frame}

\begin{frame}{Exercise 1.7.1}
  Show that if
  \[
    \int_X\int_Y |f(x,y)|\,\mu_2(dy)\,\mu_1(dx)<\infty,
  \]
  then
  \[
    \int_X\int_Y f(x,y)\,\mu_2(dy)\,\mu_1(dx)
    =
    \int_{X\times Y} f\,d(\mu_1\times\mu_2)
    =
    \int_Y\int_X f(x,y)\,\mu_1(dx)\,\mu_2(dy).
  \]
  \vikiname{Fubini and Tonelli theorem; ID: durrett_ch1_fubini_tonelli}
\end{frame}

\begin{frame}{Exercise 1.7.1: Proof Skeleton}
  \proofblock{
  Apply Tonelli to $|f|$. The hypothesis gives
  \[
    \int_{X\times Y}|f|\,d(\mu_1\times\mu_2)<\infty,
  \]
  so $f$ is integrable with respect to the product measure. Fubini now applies and gives both iterated integrals and the product integral. The same argument also justifies the common corollary
  \[
    \sum_n\int f_n\,d\Prob=\int \sum_n f_n\,d\Prob
  \]
  under absolute summability.
  }
  \vikiname{Integrable functions; ID: durrett_ch1_integrable_function}
\end{frame}

\begin{frame}{Exercise 1.7.2}
  Let $g\ge0$ be measurable on $(X,\A,\mu)$. Use Fubini/Tonelli to show
  \[
    \int_X g\,d\mu
    =
    (\mu\times\Leb)(\{(x,y):0\le y<g(x)\})
    =
    \int_0^\infty \mu(\{x:g(x)>y\})\,dy.
  \]
  \vikiname{Layer-cake formula via product sets; ID: exercise_method_layer_cake_formula}
\end{frame}

\begin{frame}{Exercise 1.7.2: Proof Skeleton}
  \proofblock{
  Let
  \[
    E=\{(x,y):0\le y<g(x)\}.
  \]
  By measurable sections, $E$ is measurable and
  \[
    \Leb(E_x)=g(x),\qquad E^y=\{x:g(x)>y\}.
  \]
  Tonelli gives
  \[
    (\mu\times\Leb)(E)=\int_X \Leb(E_x)\,\mu(dx)=\int_X g\,d\mu
  \]
  and also
  \[
    (\mu\times\Leb)(E)=\int_0^\infty \mu(E^y)\,dy.
  \]
  }
  \vikiname{Measurable sections; ID: durrett_ch1_sections_measurable}
\end{frame}

\begin{frame}[shrink=6]{Exercise 1.7.3}
  Let $F,G$ be Stieltjes measure functions and let $\mu,\nu$ be the
  corresponding measures. Show:
  \begin{enumerate}
    \item
    \[
      \int_{(a,b]}(F(y)-F(a))\,dG(y)
      =
      (\mu\times\nu)(\{(x,y):a<x\le y\le b\}).
    \]
    \item
    \[
      \int_{(a,b]}F(y)\,dG(y)+\int_{(a,b]}G(y)\,dF(y)
      =
      F(b)G(b)-F(a)G(a)
      +
      \sum_{x\in(a,b]}\mu(\{x\})\nu(\{x\}).
    \]
    \item If $F=G$ is continuous, then
    \[
      \int_{(a,b]}2F(y)\,dF(y)=F^2(b)-F^2(a).
    \]
  \end{enumerate}
  \vikiname{Stieltjes integration by parts with jump correction; ID: exercise_method_stieltjes_integration_by_parts_with_jumps}
\end{frame}

\begin{frame}{Exercise 1.7.3: Part (i)}
  \proofblock{
  Use the Stieltjes identity
  \[
    F(y)-F(a)=\mu((a,y]).
  \]
  Then
  \[
    \int_{(a,b]}(F(y)-F(a))\,\nu(dy)
    =
    \int_{(a,b]}\int \one_{\{a<x\le y\}}\,\mu(dx)\,\nu(dy).
  \]
  Tonelli turns this into the product measure of
  \[
    \{(x,y):a<x\le y\le b\}.
  \]
  }
  \vikiname{Stieltjes measure; ID: durrett_ch1_stieltjes_measure}
\end{frame}

\begin{frame}{Exercise 1.7.3: Parts (ii)--(iii)}
  \longproofblock{
  Let
  \[
    D=\{(x,y):a<x\le y\le b\},\qquad
    D'=\{(x,y):a<y\le x\le b\}.
  \]
  The sets $D$ and $D'$ cover $(a,b]^2$, and their overlap is the diagonal $\{(x,x):a<x\le b\}$. Hence
  \[
    (\mu\times\nu)(D)+(\mu\times\nu)(D')
    =
    \mu((a,b])\nu((a,b])
    +
    \sum_{x\in(a,b]}\mu(\{x\})\nu(\{x\}).
  \]
  Substitute part (i) and its symmetric version. Since $\mu((a,b])=F(b)-F(a)$ and $\nu((a,b])=G(b)-G(a)$, rearrange to obtain (ii). If $F=G$ is continuous, all point masses vanish, giving (iii).
  }
  \vikiname{Fubini and Tonelli theorem; ID: durrett_ch1_fubini_tonelli}
\end{frame}

\begin{frame}{Exercise 1.7.4}
  Let $\mu$ be a finite measure on $\R$ and
  \[
    F(x)=\mu((-\infty,x]).
  \]
  Show that, for $c>0$,
  \[
    \int_{\R}(F(x+c)-F(x))\,dx=c\,\mu(\R).
  \]
  \vikiname{Sliding interval Fubini identity; ID: exercise_method_interval_sliding_fubini}
\end{frame}

\begin{frame}{Exercise 1.7.4: Proof Skeleton}
  \proofblock{
  Write
  \[
    F(x+c)-F(x)=\mu((x,x+c])
    =
    \int_\R \one_{\{x<t\le x+c\}}\,\mu(dt).
  \]
  Tonelli gives
  \[
    \int_\R(F(x+c)-F(x))\,dx
    =
    \int_\R \Leb\{x:x<t\le x+c\}\,\mu(dt).
  \]
  For fixed $t$, the set is $[t-c,t)$, which has length $c$. Therefore the integral equals $c\mu(\R)$.
  }
  \vikiname{Fubini and Tonelli theorem; ID: durrett_ch1_fubini_tonelli}
\end{frame}

\begin{frame}{Exercise 1.7.5}
  Show that $e^{-xy}\sin x$ is integrable on the strip
  \[
    0<x<a,\qquad 0<y<\infty.
  \]
  Evaluate the double integral in both orders and obtain the usual sine-integral estimate.
  \vikiname{Sine integral via Fubini; ID: exercise_method_sine_integral_fubini}
\end{frame}

\begin{frame}{Exercise 1.7.5: Integrability}
  \proofblock{
  Use Tonelli on the absolute value:
  \[
    \int_0^a\int_0^\infty e^{-xy}|\sin x|\,dy\,dx
    =
    \int_0^a \frac{|\sin x|}{x}\,dx<\infty.
  \]
  The last integral is finite because $|\sin x|/x$ is bounded near $0$ and integrable on every compact interval away from $0$.
  }
  \vikiname{Integrable functions; ID: durrett_ch1_integrable_function}
\end{frame}

\begin{frame}[shrink=6]{Exercise 1.7.5: Two Orders}
  \longproofblock{
  Integrating in $y$ first gives
  \[
    \int_0^a\int_0^\infty e^{-xy}\sin x\,dy\,dx
    =
    \int_0^a \frac{\sin x}{x}\,dx.
  \]
  Integrating in $x$ first gives
  \[
    \int_0^a e^{-xy}\sin x\,dx
    =
    \frac{1-e^{-ay}(y\sin a+\cos a)}{1+y^2}.
  \]
  Therefore
  \[
    \int_0^a\frac{\sin x}{x}\,dx
    =
    \frac{\pi}{2}
    -
    \cos a\int_0^\infty \frac{e^{-ay}}{1+y^2}\,dy
    -
    \sin a\int_0^\infty \frac{y e^{-ay}}{1+y^2}\,dy.
  \]
  }
  \vikiname{Fubini and Tonelli theorem; ID: durrett_ch1_fubini_tonelli}
\end{frame}

\begin{frame}{Exercise 1.7.5: Estimate}
  \proofblock{
  For $a\ge1$,
  \[
    \left|\int_0^a\frac{\sin x}{x}\,dx-\frac{\pi}{2}\right|
    \le
    \int_0^\infty e^{-ay}\,dy
    +
    \int_0^\infty y e^{-ay}\,dy
    =
    \frac1a+\frac1{a^2}
    \le \frac2a.
  \]
  This is the standard quantitative form of convergence of the sine integral to $\pi/2$.
  }
  \vikiname{Sine integral via Fubini; ID: exercise_method_sine_integral_fubini}
\end{frame}

\begin{frame}{Chapter 1.7 Summary}
  \begin{itemize}
    \item Product measures are determined by measurable rectangles.
    \item Section arguments reduce product-space questions to one-coordinate questions.
    \item Tonelli handles nonnegative functions, even with infinite values.
    \item Fubini handles signed functions after absolute integrability is known.
    \item Many probability identities are just product-measure counts in disguise.
  \end{itemize}
  \vikiname{Product measures; ID: durrett_ch1_product_measure}
\end{frame}

\end{document}
